
\section{Results}
\label{sec:results}

\subsection{Identified model}
\label{sec:results:id}

Table~\ref{tab:sysid} reports the per-DOF parameters. Both translational
channels identify as fast FOPDT models, with hold-out fit $86.7\%$ ($v_x$) and $87.1\%$ ($v_y$); against the previous, slower parameters, they improve every event-by-horizon cell of the $\SEtwo$ gate. Yaw is identified from a separate open-loop dataset at $71.1\%$.

%==============================================================================
% TABLE: identified per-DOF FOPDT parameters + hold-out fit.
% Source: config/twin_params.json ; models/sysid_vxvy_v2.mat ;
%         models/sysid_omega_rotate_v1.mat ; results/identificacao_omega_20260724.md ;
%         docs/RELATORIO_TECNICO_DESENVOLVIMENTO.md sec.5 ; CHANGELOG 2026-06-14 / 2026-07-24.
%==============================================================================
\begin{table}[t]
  \caption{Identified per-DOF FOPDT parameters and hold-out fit. All three
  channels are identified (no placeholder); $\omega$ uses a dedicated
  open-loop rotation-only collection (see Sec.~\ref{sec:results:omega}), not
  the closed-loop competition logs used for $v_x,v_y$.}
  \label{tab:sysid}
  \centering
  \begin{tabular}{lcccc}
    \toprule
    DOF & $K$ & $\tau$ (s) & $\Td$ (s) & hold-out fit \\
    \midrule
    $v_x$     & 1.119 & 0.072 & 0.069 & 86.7\,\% \\
    $v_y$     & 1.052 & 0.127 & 0.056 & 87.1\,\% \\
    $\omega$  & 1.002 & 0.083 & 0.089 & 71.1\,\% \\
    \bottomrule
  \end{tabular}
  \\[2pt]
  \raggedright\footnotesize Fit is Eq.~\ref{eq:fit}, averaged over hold-out
  events. Because $v_x$ has a single hold-out event, we also ran
  leave-one-event-out over all events of each translational label: $v_x$ gives
  $84.3\pm2.0$~p.p.\ over 4 folds and $v_y$ $82.8\pm2.7$~p.p.\ over 13, with the
  refitted gain within $0.015$ of the promoted value in both cases. No parameter
  intervals are reported: see Sec.~\ref{sec:results:omega}.
  % source: results/loo_translational_20260810_1743.md — refit completo por
  %   dobra, sem alterar o modelo promovido. ATENCAO: o artefato _1742 e a
  %   execucao QUEBRADA (posicao em mm tratada como m, K~1073, tau negativo);
  %   foi descartado. Citar sempre _1743.
\end{table}


\begin{figure}[!t]
  \centerline{\includegraphics[width=0.70\columnwidth]{omega_stepresponse_en.png}}
  \caption{Measured vs.\ simulated yaw rate on the dedicated open-loop
  staircase (one training, two hold-out steps). Same model structure and
  identification tooling as on the competition logs; only the excitation regime
  differs.}
  \label{fig:fits}
\end{figure}

\subsection{Motion distribution of the hold-out}
\label{sec:results:bins}
The hold-out windows are overwhelmingly translation-dominated: of $265$ dense windows ($49$ non-overlapping), $264$ fall in the translation bin ($99.6\%$), one is rotation-dominated, and none are mixed. Re-classifying after $\omega$ was identified reproduces the same $264/1/0$ split and leaves both other bins with $N_{\mathrm{indep}}=0$, now for lack of rotation-rich competition windows rather than for lack of a yaw model.

\begin{figure*}[!t]
  \centerline{\includegraphics[width=0.86\textwidth]{main_result_en.png}}
  \caption{(a) Median free-run position error against horizon on
  translation-dominated hold-out windows, for the FOPDT model, ideal kinematics, and an untuned rSim. The two baselines cross near $\SI{1}{s}$. (b) The same   comparison resolved by event at $h_{\mathrm{ref}}$: each line is one of the   ten hold-out events, and the model is lower in all ten.
  }
  \label{fig:main}
\end{figure*}

\subsection{Translational fidelity}
\label{sec:results:rsim}

Ideal kinematics is the primary comparator: both plants run in the same
harness, start from the same measured body velocity and are scored on the same
windows, so the model is the only difference. Its median position error is
lower at every horizon (Fig.~\ref{fig:main}, Table~\ref{tab:three}):
$\SI{0.039}{m}$ against $\SI{0.101}{m}$ at $h_{\mathrm{ref}}$, and
$\SI{0.119}{m}$ against $\SI{0.291}{m}$ at $\SI{1.5}{s}$. For scale, the robot
covers about $\SI{0.4}{m}$ within $h_{\mathrm{ref}}$ and its contact diameter is
$\SI{0.16}{m}$: the baseline misses by two thirds of a robot, the model by a
quarter.

The hold-out is ten events, $\SI{73.5}{s}$ of scored motion. Windows from one
event share robot, surface, session and controller, so we take the \emph{event}
as the unit of analysis. The model wins all ten, by a median of
$+\SI{0.069}{m}$ (one-sided Wilcoxon $p = 9.8\times10^{-4}$, the smallest value
attainable at $n = 10$); the pre-registered window-level test agrees and is the
less conservative of the two. Fig.~\ref{fig:trajlabels} shows the same comparison in the trajectory domain, one panel per motion label.

%==============================================================================
% TABELA: e_pos e e_ang das TRES plantas por horizonte (bin translacional).
% Criada 2026-08-10 em resposta a B3 da revisao V2 ("o resultado principal do
% artigo nao tem figura nem tabela") e a C6 (o artigo afirmava que o e_ang do
% twin excede o do rSim em horizontes longos sem nunca dar o numero do rSim).
% Fonte: results/epos_eang_three_models_20260810_1022.md, gerada por
% scripts/epos_eang_three_models.py sobre as mesmas janelas do gate.
% NOTA: os numeros de e_ang REFUTAM a afirmacao anterior do artigo -- o modelo
% tem heading MENOR que o rSim em TODOS os horizontes. Texto corrigido.
%==============================================================================
\begin{table}[t]
  \caption{Median free-run error on translation-dominated hold-out windows.
  Position error in meters, heading error in radians. The model is lowest in
  both quantities at every horizon; the two baselines cross in position near
  $\SI{1}{s}$.}
  \label{tab:three}
  \centering
  \setlength{\tabcolsep}{4pt}
  \begin{tabular}{ccccccc}
    \toprule
    & \multicolumn{3}{c}{$\epos$ (m)} & \multicolumn{3}{c}{$\eang$ (rad)} \\
    \cmidrule(lr){2-4}\cmidrule(lr){5-7}
    $h$ (s) & model & ideal & rSim & model & ideal & rSim \\
    \midrule
    0.10 & \textbf{0.008} & 0.026 & 0.057 & \textbf{0.013} & 0.023 & 0.026 \\
    0.20 & \textbf{0.018} & 0.050 & 0.092 & \textbf{0.031} & 0.044 & 0.046 \\
    0.40 & \textbf{0.039} & 0.101 & 0.129 & \textbf{0.044} & 0.067 & 0.064 \\
    0.70 & \textbf{0.062} & 0.152 & 0.174 & \textbf{0.063} & 0.076 & 0.079 \\
    1.00 & \textbf{0.092} & 0.217 & 0.213 & \textbf{0.073} & 0.092 & 0.086 \\
    1.50 & \textbf{0.119} & 0.291 & 0.289 & \textbf{0.090} & 0.095 & 0.094 \\
    \bottomrule
  \end{tabular}
\end{table}


\paragraph{What produces the improvement}
We switch each of $\{K,\tau,\Td\}$ off in every combination and average over all
orderings, so no term is credited merely for being enabled first. Of the
$\SI{0.062}{m}$ separating the two plants at $h_{\mathrm{ref}}$, the time
constant accounts for $55\%$, the dead time for $43\%$ and the gain for $2\%$.
Lag and delay are comparable contributors; the gain is irrelevant.

\paragraph{The improvement needs the measured initial velocity}
Both plants are seeded with the measured body velocity, but only the dynamic one
can carry it forward, and neither ingredient works alone. Extrapolating that
velocity at constant speed, ignoring the commands, gives $\SI{0.160}{m}$;
running the model from rest gives $\SI{0.114}{m}$. Both are worse than ideal
kinematics, and only the combination reaches $\SI{0.039}{m}$. We treat the
measured velocity as part of the model interface rather than privileged
information, since a twin running alongside the robot has that estimate;
re-deriving it with a strictly causal difference raises both errors but
\emph{widens} the gap, so the verdict does not rest on it.

\paragraph{The ODE simulator}
rSim reaches $\SI{0.129}{m}$ at $h_{\mathrm{ref}}$, so even the zero-parameter baseline outperforms it. Giving it the model's dead time widens rather than closes the gap.

\paragraph{Heading error}
The rotation and mixed bins lack a scorable population
(Sec.~\ref{sec:results:bins}), so heading carries no verdict. Descriptively
(Table~\ref{tab:three}), the model's median $\eang$ is the lowest of the three
plants at every horizon, but the margin shrinks from roughly half the
baselines' at $\SI{0.1}{s}$ to within $\SI{0.005}{rad}$ of both by
$\SI{1.5}{s}$. The position advantage does not transfer to heading in the same
proportion, which is consistent with yaw being the channel the competition logs
cannot constrain.

\subsection{Yaw identifiability is a function of excitation regime}
\label{sec:results:omega}

\begin{figure*}[!t]
  \centerline{\includegraphics[width=0.98\textwidth]{trajectory_by_label_en.png}}
  \caption{Free-run trajectory overlays, one per motion label, over a
  $\SI{1.5}{s}$ window re-anchored at the black dot. For each label we show the
  window with the largest margin of the model over the better of the two
  comparison plants, so these are per-label best cases rather than typical
  windows; the median-window overlay and the full distributions are in
  Fig.~\ref{fig:main} and Table~\ref{tab:three}. Panel titles give the
  end-of-window position error of each plant.}
  \label{fig:trajlabels}
\end{figure*}

The problem is visible in the logs before any model is fitted. Regressing the
measured body velocity on the command, aligned by its dead time, the command
explains $96\%$ of the variance for $v_x$ and $94\%$ for $v_y$, but only $56\%$
for $\omega$: nearly half the rotation the cameras see in play is not
attributable to the yaw command.

Fitting the same structure to the highest-rotation competition events gives
$K_\omega = 1.378$ and a hold-out fit of $44.8\%$, below the pre-registered
$60\%$ gate. Using it also \emph{worsens} free-run position, because a gain
above unity amplifies accumulated heading drift, so it is not promoted.

\paragraph{Why the logs cannot constrain yaw}
The firmware clamps each \emph{wheel} setpoint below $\SI{1.567}{rad/s}$, after
the command kinematics. Converted to body yaw rate that threshold sits at
$0.58$--$0.60\,\si{rad/s}$, and $94.3\%$ of commanded yaw samples in the logs
fall below it. The statistic is easy to misread: it does not mean yaw is
switched off. In play the wheels are already turning because of translation, so
among those samples the median per-wheel translational setpoint is $15\times$
the clamp and only $2.2\%$ are fully clamped. The yaw command is not suppressed
but \emph{diluted}, contributing a median of $1.0\%$ of what each wheel is asked
to do. What the actuators receive in competition is $99\%$ translation.

\paragraph{The same structure, different data}
A short dedicated open-loop collection removes the dilution: with
$v_x = v_y = 0$ the yaw command becomes the sole actuation. Identified with the
\emph{same} structure and tooling (Fig.~\ref{fig:fits}), it gives
$K_\omega = 1.002$, $\tau_\omega = \SI{0.083}{s}$, $\Td{}_\omega = \SI{0.089}{s}$
and a hold-out fit of $71.1\%$, above the gate. Scored against the same
placeholder on the same competition population, the closed-loop estimate is
worse at every horizon and the open-loop one better, so the verdict inverts
under identical tooling.